\input{\string~/Documents/School/"UC Irvine"/ta_template2.0.tex}
\rh{2B}{7.4}
\chead{\today}
\graphicspath{ {\string~/Desktop} }
\usetikzlibrary{arrows}
%============ANSWER KEY TOGGLE===========
\newcommand{\ans}[1]{ \noindent {\color{red} \textbf{Answer:} #1}}
\renewcommand{\ans}[1]{\vspace{3in}}
%==========================================
\begin{document}
\begin{center}
\textbf{Partial Fraction Decomposition}
\end{center}
\begin{aun}
In partial fractions, we have a rational function:
\[
\frac{P(x)}{Q(x)} = S(x)+ \frac{R(x)}{Q(x)}
\]
after long division. Note that, once we complete this, the degree of $R$ is less than the degree of $Q$. Integrating the second rational polynomial, we can break into a few cases based on what's going on with the denominator, $Q(x)$:
\begin{description}
\item[$Q(x)$ factors into distinct \underline{linear} factors:] We can write:
\[
Q(x)=\left(a_1 x+b_1\right)\left(a_2 x+b_2\right) \cdots\left(a_k x+b_k\right) \mte{ and } \frac{R(x)}{Q(x)}=\frac{A_1}{a_1 x+b_1}+\frac{A_2}{a_2 x+b_2}+\cdots+\frac{A_k}{a_k x+b_k}
\]
\item[$Q(x)$ has some repeated linear factors:] In this case, we need to account for the repetitions by adding extra fraction terms:
\[
Q(x)=\left(a_1 x+b_1\right)^r \mte{ and }  \frac{R(x)}{Q(x)} = \frac{A_1}{a_1 x+b_1}+\frac{A_2}{\left(a_1 x+b_1\right)^2}+\cdots+\frac{A_r}{\left(a_1 x+b_1\right)^r} + \cdots 
\]
\item[$Q(x)$ has some \underline{quadratic} factors:] In this case, our numerator will not only have constants, but needs to have an entire \emph{linear} function:
\[
Q(x) = \underbrace{ax^2 + bx + c}_{b^2 - 4ac < 0} \cdots (\text{other factors})\mte{ and } \frac{R(x)}{Q(x)} =  \frac{A x+B}{a x^2+b x+c} + \cdots
\]
\item[$Q(x)$ has repeated quadratic factors:] In this case, we just combine the techniques from the preceding two cases, where we have repeated linear terms for each time the quadratic is repeated:
\[
Q(x) = \underbrace{(ax^2 + bx + c)^r}_{b^2 - 4ac < 0} \cdots (\text{other factors}) \]
and  
\[
\frac{R(x)}{Q(x)} = \frac{A_1x + B_1}{ax^2 + bx + c}+\frac{A_2x + B_2}{\left(ax^2 + bx + c\right)^2}+\cdots+\frac{A_rx + B_r}{\left(ax^2 + bx + c\right)^r} + \cdots 
\]
\end{description}
\label{pfd}
\end{aun}
\begin{exx}
Compute the following integral:
\[
\int \frac{4+x}{(1+2 x)(3-x)} dx
\]
\end{exx}
\ans{Since we're in the case of distinct linear factors, we'll set up:
\[
\frac{4+x}{(1+2 x)(3-x)} = \frac{A}{2x + 1} + \frac{B}{3 - x}
\]
Clearing denominators gives us:
\[
4 + x = A (3 - x) + B(2x + 1)
\]
Plugging in $x = 3$, we see:
\[
4 + 3 = \cancelto{0}{A (3 - 3)} + B(2(3) + 1) \mte{ implying that } B = 1
\]
Plugging in $x = -\frac12$ gives us:
\[
4 - \frac12 = A (3 + \frac12) + \cancelto{0}{B (2 (-\frac12) + 1)} \mte{ implying } A = 1
\]
This allows us to simplify our integrand into two integrals:
\[
\int \frac{4+x}{(1+2 x)(3-x)} dx = \int \frac{dx}{2x + 1} + \int \frac{dx}{3 - x}
\]
}

\problembox{Compute the following:
\[
\int \frac{x^3+4 x+3}{x^4+5 x^2+4} d x
\]}
\ans{
Denominator factors into:
\[
\frac{x^3 + 4x + 3}{(x^2 + 4)(x^2 + 1)} = \frac{Ax + B}{x^2 + 4} + \frac{Cx + D}{x^2 + 1}
\]
We can solve for $A,B,C, D$ algebraically:
\[
(Ax+ B)(x^2 + 1) + (Cx + D)(x^2 + 4) = Ax^3 + Ax + Bx^2 + B + Cx^3 + 4 Cx + Dx^2 + 4D = x^3 + 0x^2 + 4x + 3 
\]
From here, just pair off our coefficients based on degree. The terms on both sides have:
\[
\begin{array}{|c|c|}
\hline
\textbf{Term} & \textbf{Coefficients} \\ \hline
x^3 & A + C = 1\\
x &  A + 4C = 4\\
x^2 & B + D = 0 \\
c & B + 4D = 3 \\ \hline
\end{array}
\]
This is a system of equations in 4 variables, with the solution $A = 0$, $B = -1$, $C = 1$, $D = 1$, completing our decomposition. From here, we may rewrite the integrand as:
\[
\int \frac{-1}{x^2 + 4} + \frac{x + 1}{x^2 + 1} dx
\]
which we can integrate. 
}

\problembox{[Additional Practice] Consider the following integrals:
\begin{enumerate}
\item $\dst \int \frac{x^2-3 x+7}{\left(x^2-4 x+6\right)^2} d x$
\item $\dst \int \frac{\sec ^2 t}{\tan ^2 t+3 \tan t+2} d t$
\end{enumerate}}
\ans{
\begin{enumerate}
\item A denominator like this is going to have repeated factors, so we need to remember to address this in our solution. 
\item Use a $u$ substitution for $u = \tan t$ and then this condenses into a computation very similar to the first example. 
\end{enumerate}
}
\end{document}
